% Partie : théorie de la démonstration
% Chapitre : élimination des coupures
% Section : introduction

\documentclass[../../../include/open-logic-section]{subfiles}

\begin{document}

\olfileid[fr]{pt}{cut}{int}

\olsection{Introduction}

La règle \Cut{} est la suivante :
\[
\Axiom$\Gamma \fCenter \Delta, !A$
\Axiom$!A, \Pi \fCenter \Lambda$
\RightLabel{\Cut}
\BinaryInf$\Gamma, \Pi \fCenter \Delta, \Lambda$
\DisplayProof
\]
La !!{formula}~$!A$ est appelée \emph{formule de coupure}.

La règle \Cut{} figurait dans le calcul des séquents~\Log{LK}
original de Gentzen. Son \emph{Hauptsatz}, ou « théorème principal »,
affirme que tout séquent !!{provable} dans~\Log{LK} est aussi
!!{provable} sans la règle \Cut{} : les coupures peuvent être
éliminées des preuves de~\Log{LK}. Nos systèmes \Log{G1c} et
\Log{G3c} ne contiennent pas \Cut, mais la même question se pose :
peut-on éliminer \Cut{} des !!{proof}s qui emploient cette règle
en plus de celles de \Log{G1c}, ou de \Log{G3c} ?
Dans la terminologie déjà introduite, ceci revient à demander
si \Cut{} est admissible dans \Log{G1c}, ou dans \Log{G3c}.

Le théorème d'élimination des coupures est peut-être le résultat
le plus important et le plus célèbre de la théorie de la démonstration.
Il établit d'abord que \Log{G1c} et \Log{G3c} se passent d'une règle
qui semble, à première vue, indispensable. En formalisant des
démonstrations usuelles, on doit en effet pouvoir employer des lemmes :
les mathématiciens démontrent des résultats, puis s'en servent pour
en démontrer d'autres. On peut formaliser la preuve d'un résultat $L$,
le lemme, par une !!{proof} du séquent $\Sequent L$.
On formalise ensuite la preuve du second résultat~$R$, qui emploie~$L$,
par une !!{proof} du séquent $L \Sequent R$.
Comment savoir qu'il existe une preuve de~$R$, c'est-à-dire une
!!{proof} du seul séquent $\Sequent R$ ?
Il faut pouvoir combiner les deux !!{proof}s ; la règle \Cut{}
le permettrait.

Nous avons indiqué que \Log{G1c} et \Log{G3c} sont complets :
ils démontrent tous les séquents valides. Ce résultat peut être
établi directement. À l'époque où Gentzen développait ses travaux,
seule la complétude de systèmes axiomatiques de dérivation était
connue. Pour établir celle de \Log{LK}, il a traduit les
!!{derivation}s du système axiomatique en !!{proof}s de~\Log{LK}.
Cette traduction faisait un usage essentiel de \Cut{}.

L'élimination des coupures concerne aussi d'autres logiques,
qui possèdent elles aussi des calculs de séquents.
Pour certaines, une preuve directe de complétude peut être difficile,
voire hors de portée des méthodes disponibles. Une traduction
depuis un autre système, faisant intervenir \Cut, fournit alors
une voie pour établir la complétude. Une preuve de l'élimination
des coupures par les méthodes de la théorie de la démonstration
permet ensuite d'établir que \Cut{} est superflu et que le système
sans \Cut{} est complet.

L'élimination des coupures permet également d'obtenir d'autres
résultats présentant un intérêt propre. Nous examinerons deux
applications : le lemme d'interpolation de Craig et le théorème
de Herbrand.

Les preuves de l'élimination des coupures montrent généralement
comment transformer une !!{proof} utilisant \Cut{} en une preuve
sans coupure. Ces transformations sont habituellement locales :
on remplace une portion d'une !!{proof}, généralement une
sous-!!{proof} se terminant par \Cut, par une autre sous-!!{proof}
du même séquent, où cette coupure a été supprimée ou remplacée
par d'autres. Ces arguments donnent des algorithmes qui transforment,
étape par étape, les !!{proof}s avec coupures en preuves sans coupure.

Cette information supplémentaire est utile dans les applications.
Par exemple, une preuve par la théorie des modèles du lemme
d'interpolation établit que, si $!A \lif !B$ est valide, un
interpolant existe ; laissons pour le moment sa définition de côté.
Un argument fondé sur l'élimination des coupures donne un algorithme
qui prend une !!{proof} de $!A \lif !B$ et produit un interpolant.
Cet argument de théorie de la démonstration est donc
\emph{constructif}.

Deux méthodes usuelles permettent de démontrer l'élimination des
coupures. La première, qui remonte à Gentzen, supprime les coupures
les plus hautes d'une preuve et répète l'opération jusqu'à ce
qu'il n'en reste aucune. La seconde, due à Sch\"utte et Tait,
vise les coupures les plus complexes : on supprime successivement
les coupures maximales, au sens où aucune autre coupure ne porte
sur une !!{formula} de profondeur supérieure.
Chaque méthode a ses avantages et ses limites. Pour certains
systèmes, l'une se met aisément en œuvre tandis que l'autre
est difficile, voire ne s'applique pas sous la forme envisagée.
Nous présenterons donc les deux.

Nous démontrerons l'élimination des coupures pour~\Log{G3c}.
Plus précisément, nous éliminerons la \emph{coupure à contexte commun}
\CutCS{} :
\[
\Axiom$\Gamma \fCenter \Delta, !A$
\Axiom$!A, \Gamma \fCenter \Delta$
\RightLabel{\CutCS}
\BinaryInf$\Gamma \fCenter \Delta$
\DisplayProof
\]
Si un calcul des séquents autorise l'affaiblissement et la contraction,
comme règles primitives dans \Log{G1c} ou comme règles admissibles
dans \Log{G3c}, \Cut{} permet de simuler \CutCS, et réciproquement.
Nous choisissons \CutCS{} parce que la première méthode est plus
simple à décrire sous cette forme, et que la seconde méthode sera
présentée avec cette version à contexte commun.











\end{document}
